\documentclass{internal-report}
\usepackage[UTF8]{ctex}
\usepackage{booktabs}
\usepackage{amsmath,amssymb}

% STATUS: draft
% LAST_MODIFIED: 2026-08-09

\title{S1 预测段佐证：GPU 需求序列不可预测性的四重统计检验}
\author{MathModel 建模小组}
\date{2026-08-09}

\begin{document}

\maketitle

\begin{abstract}
问题一要求对不同区域、不同类型任务的 GPU 需求做统计分析并构建短期预测模型。
阶段 1.1 的 A 类验证（\texttt{verify-sub1-20260807.md}）已通过 ACF 诊断与泊松
拟合给出"逐时 GPU 需求序列近乎白噪声"的初步结论。本文件对该结论做\textbf{四重
对抗性检验}（B 类 head-to-head）：(T1) BDS 非线性依赖检验——排除任意维度的
（线性与非线性）时序依赖；(T2) Prophet 对照——排除可被傅里叶分解捕捉的周期
结构；(T3) LSTM 对照——排除可被循环神经网络学习的非线性记忆；(T4) 跨类型/
跨区域 Granger 因果——排除可利用的横截面外生路径。四项检验均无法推翻白噪声
结论，且最复杂模型（Prophet）相对常数均值基线的测试窗提升仅 $+2.9\%$（小于
$5\%$ 判据），故预测段采用"统计检验完备的白噪声论证 + 简单稳健基线（常数均值
/线性回归）"为最终叙事。本文件为论文问题一预测部分的公式化佐证章节。
\end{abstract}

\section{预测目标与赛题口径}

\subsection{预测目标序列}

设第 $t$ 小时到达任务集合为 $\mathcal{J}_t=\{j: A_j=t\}$，则逐时 GPU 需求序列

\begin{equation}
  Y_t = \sum_{j \in \mathcal{J}_t} g_j,\qquad t=0,1,\ldots,2399,
  \label{eq:yseq}
\end{equation}

其中 $A_j$、$g_j$ 分别为任务 $j$ 的到达小时与 GPU 需求。分类型序列
$Y_t^{(k)}$（$k\in\{\text{训练},\text{批量},\text{实时}\}$）定义同理。

\subsection{赛题协议切分}

赛题说明 5 指定三段式协议，全程遵守（\textbf{无测试窗信息泄漏}）：

\begin{center}
\begin{tabular}{lll}
\toprule
阶段 & 小时区间 & 用途 \\
\midrule
训练 & $0$--$2351$ & 估计模型参数 \\
调参验证 & $2352$--$2375$ & 超参数/早停选择 \\
重训 & $0$--$2375$ & 最终训练（含调参窗） \\
测试 & $2376$--$2399$ & 唯一评价窗口 \\
\bottomrule
\end{tabular}
\end{center}

\subsection{基准模型（诚实口径）}

\begin{itemize}
  \item 常数均值：$\hat{Y}_t = \bar{Y}_{0:2376} = \frac{1}{2376}\sum_{t=0}^{2375} Y_t$（训练窗均值外推）。
  \item 线性回归：$\hat{Y}_t = \beta_0 + \beta_1 t + \sum_{p\in\{24,168\}}\bigl(\alpha_p\sin\frac{2\pi t}{p} + \gamma_p\cos\frac{2\pi t}{p}\bigr)$，系数在训练窗最小二乘估计。
\end{itemize}

\section{T1\quad BDS 非线性依赖检验}

\subsection{检验原理}

ACF $\hat{\rho}_k \approx 0$ 只能排除\textbf{线性}自相关；若序列存在非线性
依赖（如 GARCH 类波动聚集、阈值效应），LSTM 等模型仍可能从中获益。BDS 检验
（Brock--Dechert--Scheinkman, 1996）以嵌入维度 $m$ 上的相关积分 $C_m(\varepsilon)$
为统计量，对任意 $m$ 维依赖结构敏感，原假设为 $H_0$：序列 $i.i.d.$。

对样本 $\{Y_t\}_{t=1}^{T}$，定义 $m$ 维嵌入向量 $Y_t^m=(Y_t,Y_{t-1},\ldots,Y_{t-m+1})$，
其相关积分

\begin{equation}
  C_m(\varepsilon) = \frac{2}{(T-m+1)(T-m)}
    \sum_{s<t} \prod_{i=0}^{m-1} \mathbf{1}\bigl\{|Y_{s+i}-Y_{t+i}| \le \varepsilon\bigr\},
  \label{eq:corrintegral}
\end{equation}

$\varepsilon$ 取 $\varepsilon = 1.5\,\hat{\sigma}_Y$（$\hat{\sigma}_Y$ 为样本标准差）。
BDS 统计量

\begin{equation}
  W_{m,T}(\varepsilon) = \sqrt{T-m+1}\,
    \frac{C_m(\varepsilon) - C_1(\varepsilon)^{m}}{\hat{\sigma}_m(\varepsilon)}
  \;\xrightarrow[H_0]{d}\;\mathcal{N}(0,1),
  \label{eq:bdsstat}
\end{equation}

其中 $\hat{\sigma}_m^2(\varepsilon)$ 为 $C_m(\varepsilon)$ 在原假设下的渐近方差
（由 $m$ 个相关积分 $C_1,\ldots,C_m$ 的显式组合给出，见 Brock et al., 1996）。
双侧 $p<0.05$ 拒绝 $H_0$，判定存在依赖。

\subsection{检验结果}

\begin{center}
\footnotesize
\begin{tabular}{lcccc}
\toprule
序列 & $p$（$m=2$） & $p$（$m=3$） & $p$（$m=4$） & 判定 \\
\midrule
Total（逐时 GPU 需求） & 0.4004 & 0.8208 & 0.7652 & 不拒绝 $H_0$（$\approx$ i.i.d.） \\
AITraining & 0.7363 & 0.9130 & 0.7801 & 不拒绝 $H_0$ \\
BatchInference & 0.3073 & 0.5551 & 0.3440 & 不拒绝 $H_0$ \\
RealTimeInference & 0.5504 & 0.5769 & 0.7321 & 不拒绝 $H_0$ \\
\bottomrule
\end{tabular}
\end{center}

全部 $p$ 值 $\ge 0.30$，远高于 $0.05$ 显著性水平：\textbf{不仅线性自相关，
任意 $m\le 4$ 维的非线性依赖结构亦不存在}。LSTM/GARCH 类模型在此数据上无结构可学。

\section{T2\quad Prophet 对照检验}

\subsection{检验原理}

Prophet 将序列分解为趋势与傅里叶周期项：

\begin{equation}
  Y_t = g(t) + s_{\text{day}}(t) + s_{\text{week}}(t) + \varepsilon_t,
  \qquad s_p(t) = \sum_{n=1}^{N_p}\Bigl(a_n\cos\frac{2\pi n t}{p} + b_n\sin\frac{2\pi n t}{p}\Bigr),
  \label{eq:prophet}
\end{equation}

开启日周期（$p=24$）与周周期（$p=168$）傅里叶项。若序列存在 ACF 无法
捕捉但傅里叶基可表达的真实周期，Prophet 应显著优于常数均值。

\subsection{检验结果}

\begin{center}
\begin{tabular}{lcc}
\toprule
阶段 & RMSE & MAPE \\
\midrule
调参窗 $2352$--$2375$ & 186.4 & 32.7\% \\
测试窗 $2376$--$2399$ & \textbf{214.7} & 22.6\% \\
\bottomrule
\end{tabular}
\end{center}

测试窗 Prophet RMSE $=214.7$ 对常数均值 $221.1$，提升 $+2.9\% < 5\%$：
\textbf{周期结构不存在，Prophet 的分解优势无的放矢}。调参窗 RMSE（186.4）
显著低于测试窗（214.7），说明该"提升"为样本间波动，不具备泛化性。

\section{T3\quad LSTM 对照检验}

\subsection{检验原理}

采用单层 LSTM（16 隐单元，输入=滞后 24h 窗口 $[Y_{t-24},\ldots,Y_{t-1}]$，
输出=$\hat{Y}_t$，滚动一步预测）。训练于 $0$--$2351$，在调参窗 $2352$--$2375$
早停（选验证 RMSE 最优的 epoch），归一化统计量仅由训练窗估计。若存在非线性
时序记忆，LSTM 应优于均值基线。

\subsection{检验结果}

\begin{center}
\begin{tabular}{lcc}
\toprule
阶段 & RMSE & 说明 \\
\midrule
调参窗最优（早停） & 180.7 & 训练窗内表现 \\
测试窗 $2376$--$2399$ & \textbf{221.5} & 泛化表现 \\
\bottomrule
\end{tabular}
\end{center}

测试窗 LSTM RMSE $=221.5$，较常数均值 $221.1$ \textbf{更差}（$-0.2\%$）：
调参窗 180.7 到测试窗 221.5 的严重退化是典型的"对白噪声序列过拟合"。
LSTM 学到的"记忆"在测试窗无任何预测价值。

\section{T4\quad 跨类型 / 跨区域 Granger 因果检验}

\subsection{检验原理}

若某序列（如训练任务到达）能 Granger 原因解释目标序列，则可借外生路径提升
预测。对候选原因 $X$ 与目标 $Y$，估计受限与不受限两式：

\begin{equation}
  \text{受限：}\quad Y_t = \alpha_0 + \sum_{i=1}^{p}\alpha_i Y_{t-i} + u_t,
  \qquad
  \text{不受限：}\quad Y_t = \alpha_0 + \sum_{i=1}^{p}\alpha_i Y_{t-i}
    + \sum_{i=1}^{p}\beta_i X_{t-i} + v_t,
  \label{eq:granger}
\end{equation}

$H_0$：$X$ 不 Granger 原因 $Y$，即 $\beta_1=\cdots=\beta_p=0$。检验统计量

\begin{equation}
  F = \frac{(\text{RSS}_r - \text{RSS}_u)/p}{\text{RSS}_u/(T-2p-1)}
  \;\sim\; F(p,\, T-2p-1),
  \label{eq:grangerf}
\end{equation}

$p=3$，$p<0.05$ 拒绝 $H_0$。判定规则：$\ge 2$ 个滞后阶显著才判"存在因果"
（保守，抑制多重检验假阳性）。

\subsection{检验结果}

\begin{center}
\footnotesize
\begin{tabular}{lccccl}
\toprule
原因 $X$ & 目标 $Y$ & lag$=1$ & lag$=2$ & lag$=3$ & 判定 \\
\midrule
AITraining & BatchInference & 0.4013 & 0.5668 & 0.6724 & 无 \\
AITraining & RealTimeInference & 0.4697 & 0.4161 & 0.4635 & 无 \\
BatchInference & AITraining & 0.4507 & 0.6245 & 0.6576 & 无 \\
BatchInference & RealTimeInference & 0.4565 & 0.6965 & 0.8163 & 无 \\
RealTimeInference & AITraining & 0.1287 & 0.0699 & 0.1255 & 无 \\
RealTimeInference & BatchInference & 0.0607 & 0.1541 & 0.2507 & 无 \\
\midrule
RegionA & RegionB & 0.1909 & 0.2227 & 0.3928 & 无 \\
RegionA & RegionD & 0.6930 & 0.7993 & 0.6960 & 无 \\
RegionB & RegionC & 0.4725 & 0.6015 & 0.7596 & 无 \\
RegionE & RegionF & 0.7044 & 0.6903 & 0.8560 & 无 \\
\bottomrule
\end{tabular}
\end{center}

6 组类型间 + 4 组区域间配对的最小 $p$ 值 $=0.0607$（RealTimeInference
$\to$ BatchInference，lag$=1$），仍高于 $0.05$：\textbf{不存在可利用的跨序列
因果路径}，外生变量预测路线不可行。

\section{综合结论与对比表}

\subsection{四模型测试窗对比（2376--2399）}

\begin{center}
\begin{tabular}{lccc}
\toprule
模型 & RMSE & MAPE & 相对常数均值 \\
\midrule
Prophet（日/周季节） & \textbf{214.7} & 22.6\% & $+2.9\%$ \\
线性回归（周期特征） & 216.1 & 23.1\% & $+2.3\%$ \\
常数均值（基准） & 221.1 & 23.4\% & --- \\
LSTM（滚动一步） & 221.5 & 23.4\% & $-0.2\%$ \\
\bottomrule
\end{tabular}
\end{center}

\subsection{四重检验结论汇总}

\begin{center}
\footnotesize
\begin{tabular}{p{2.6cm}p{4.6cm}p{2.2cm}p{4.6cm}}
\toprule
检验 & 检验内容（$H_0$） & 结果 & 对不可预测性结论的支持 \\
\midrule
ACF（阶段 1.1） & 线性自相关 $\rho_k=0$ & 全部 $\approx 0$ & 线性依赖不存在 \\
T1 BDS & 序列 i.i.d. & 全部 $p\ge 0.30$ & 任意 $m\le4$ 维依赖不存在 \\
T2 Prophet & 傅里叶周期结构无信息 & 提升 $+2.9\%$ & 周期结构不存在 \\
T3 LSTM & 非线性记忆可学习 & 提升 $-0.2\%$ & 无非线性结构可学 \\
T4 Granger & 跨序列无因果 & 全部 $p>0.05$ & 无外生路径可利用 \\
\bottomrule
\end{tabular}
\end{center}

\subsection{结论与论文引用口径}

上述四项检验覆盖了"可预测性"的四条可能路径——线性依赖（ACF）、非线性依赖
（BDS）、周期结构（Prophet）、横截面因果（Granger）——全部未推翻白噪声结论。
即便允许最复杂模型（Prophet），其对常数均值基线的测试窗提升亦仅 $+2.9\%$，
低于论文采用的 $5\%$ 实质改善判据，且调参窗/测试窗性能倒挂（186.4 vs 214.7）
表明该差距为样本噪声而非真实信号。

因此问题一预测段采取以下口径：
\begin{enumerate}
  \item \textbf{统计发现}：逐时 GPU 需求序列为（近似）i.i.d. 白噪声，到达过程近似泊松（均值 $\approx 20.8$/h）——以 ACF + BDS + Prophet + LSTM + Granger 五组证据闭环；
  \item \textbf{预测模型}：常数均值（训练窗均值外推）为主预测，线性回归作对照，二者在测试窗 RMSE 221.1 / 216.1、MAPE 23.4\% / 23.1\%；
  \item \textbf{口径声明}：赛题说明 2 明确"第 2376--2399 小时预测精度仅用于预测评价"，末 24h 调度以实际到达任务为输入，预测误差不影响调度方案本身；
  \item \textbf{边界说明}：本文证明的是"点预测（逐时期望值）不可行"；若需概率化输出，可用泊松分布刻画到达过程（分布预测路径，见后续工作），但不改变点预测结论。
\end{enumerate}

\end{document}
